% !PRINTED 520
\weberpage{520}
%<<CONT
\weq{20}{\begin{aligned}
M &= (-1)^{\frac{m}{8}} \frac{1+i}{\sqrt2} \left(1 - i\tfrac12\sqrt{m}\right)^3, & A &\equiv 1\ (\mathrm{mod}\ 8), \\
M &= (-1)^{\frac{m}{8}} \frac{-1+i}{\sqrt2} \left(1 - i\tfrac12\sqrt{m}\right)^3 = M', & A &\equiv 3\ (\mathrm{mod}\ 8), \\
M &= (-1)^{\frac{m}{8}} \frac{-1-i}{\sqrt2} \left(1 - i\tfrac12\sqrt{m}\right)^3 = M'', & A &\equiv 5\ (\mathrm{mod}\ 8), \\
M &= (-1)^{\frac{m}{8}} \frac{1-i}{\sqrt2} \left(1 - i\tfrac12\sqrt{m}\right)^3 = M''', & A &\equiv 7\ (\mathrm{mod}\ 8).
\end{aligned}}
Die vier Werte von $M$ gehen aus dem ersten hervor durch die Vertauschungen:
\weq{21}{\begin{aligned}
M, &\quad \sqrt2, \quad i, \quad \sqrt{m}, \\
M', &\quad -\sqrt2, \quad -i, \quad -\sqrt{m}, \\
M'', &\quad -\sqrt2, \quad i, \quad \sqrt{m}, \\
M''', &\quad \sqrt2, \quad -i, \quad -\sqrt{m}.
\end{aligned}}
Ist nun $\Phi(M,u) = 0$ in demselben Sinne wie oben die zwischen $u$ und $M$ bestehende Gleichung, so sind die vier Funktionen
\[
\Phi(M,u), \quad \Phi(-M,u), \quad \Phi(iM,u), \quad \Phi(-iM,u)
\]
ohne gemeinsamen Teiler, und wenn $\Phi(M,u)$ in $H(u)$ enthalten ist, so sind auch die drei anderen Funktionen $\Phi(-M,u)$, $\Phi(iM,u)$, $\Phi(-iM,u)$ in $H(u)$ enthalten, da die Vorzeichen der Irrationalitäten (20) beliebig geändert werden können. Es ist daher
\[
H(u) = \Phi(M,u)\,\Phi(-M,u)\,\Phi(iM,u)\,\Phi(-iM,u),
\]
und es kommt also jeder der vier Fälle
\[
A \equiv 1, \quad A \equiv 3, \quad A \equiv 5, \quad A \equiv 7\ (\mathrm{mod}\ 8)
\]
in gleich vielen Formenklassen der Diskriminante $-4m$ vor.

Eine Ausnahme tritt nur dann ein, wenn $m$ ein Quadrat oder das Doppelte eines Quadrates ist, weil im ersten Falle $\sqrt m$ rational und daher nur die Vertauschung $(M, M'')$ gestattet ist, im zweiten Falle $\sqrt2$ und $\sqrt m$ gleichzeitig ihr Zeichen ändern, also nur die Vertauschung $(M, M')$ zulässig ist.

Aber es genügen auch schon diese Vertauschungen, um die Teilgleichung durch Adjunktion von $\sqrt2$ weiter zu zerfällen.
